\chapter{\textcolor{gray}{Du système nerveux au neurone}}
\section{Le système nerveux \cite{widmaier2008vanders} :}
	\begin{figure}[htb]
		\centering
		\includegraphics[width=0.6\textwidth]{pictures/systeme nerveu.png}
		\caption{Système nerveux}
	\end{figure}

	\textbf{Le système nerveux central (SNC)} comprend le cerveau et la moelle épinière .  
	Le cerveau contrôle les fonctions comme la pensée, les mouvements et les sensations.  
	La moelle épinière relie le cerveau au reste du corps et fait circuler les messages nerveux.\\[1ex]
	\textbf{Le système nerveux périphérique (SNP)} regroupe les nerfs situés en dehors du SNC.  
	Il transmet les informations entre le corps et le cerveau.  
	Il se divise en deux parties :
	\begin{itemize}
		\item \textbf{Système somatique} : contrôle les mouvements volontaires (comme marcher).
		\item \textbf{Système autonome} : gère les fonctions involontaires (comme respirer).
		Il comporte deux parties : 
		\begin{itemize}
			\item \textbf{Le système sympathique} : prépare le corps au stress ou au danger (cœur qui bat vite, respiration rapide, pupilles dilatées).
			\item \textbf{Le système parasympathique} : aide le corps à se détendre et à revenir au calme (ralentit le cœur, la respiration, et économise l'énergie).
		\end{itemize}
	\end{itemize}


\section{Le cerveau}

	\hspace*{1em}Le cerveau est l'organe situé dans la boîte crânienne, responsable de la régulation de la plupart des fonctions du corps. 
	Il traite les informations sensorielles, contrôle les mouvements, et gère les fonctions cognitives telles que la pensée, la mémoire, et les émotions.
	
	\vspace{1em}
	
	\hspace*{1em}\textbf{Le cortex cérébral} : Divisé en quatre lobes principaux (frontal, pariétal, temporal, et occipital), 
	il est responsable de fonctions complexes comme la pensée, la perception, la planification et la compréhension du langage.
	
	\vspace{0.5em}
	
	\hspace*{1em}\textbf{Le cervelet} : Situé sous le cerveau, il joue un rôle crucial dans la coordination des mouvements et l'équilibre.
	
	\vspace{0.5em}
	
	\hspace*{1em}\textbf{Le tronc cérébral} : Il contrôle les fonctions vitales telles que la respiration, la fréquence cardiaque et la digestion.    
\begin{figure}[b]
	\centering
	\includegraphics[width=0.7\linewidth]{pictures/Les différents lobes du cerveau.png} 
	\caption{Les différents lobes du cerveau}
\end{figure}
	   
\begin{figure}[h]
	\centering
	\includegraphics[width=0.9\linewidth]{pictures/l anatomie de cerveau 1.png}
	\caption{Ventricules du cerveau}
\end{figure}

\begin{figure}[H]
	\centering
	\includegraphics[width=0.9\linewidth]{pictures/les neurons dans le cerveau.png}
	\caption{Section frontale des hémisphères cérébrales, du thalamus et de l'hypothalamus}
\end{figure}
\begin{figure}[h]
		\centering
		\includegraphics[width=0.9\linewidth]{pictures/les neurons dans le cerveau 2.png} 
		\caption{Cellules gliales du système nerveux central}
\end{figure}
	
\section{Les neurones \cite{gerstner2014neuronal} \cite{kandel2012principles}}
	\begin{figure}[h]
		\centering
		\includegraphics[scale=0.2]{pictures/anatomie de neuron 2.png}
		\caption{Diagramme représentatif du neurone}
	\end{figure}

	\hspace*{1em} \textbf{Un neurone} est une cellule du système nerveux spécialisée dans l'initiation, 
	l'intégration et la conduction de signaux électriques vers d'autres cellules, parfois sur de longues distances. Un signal peut déclencher de nouveaux signaux 
	électriques dans d'autres neurones, ou stimuler une cellule glandulaire à sécréter des substances, ou encore une cellule musculaire à se contracter. 
	
	\hspace*{1em}Ainsi, les neurones constituent un moyen essentiel de contrôle des activités d'autres cellules. L'incroyable complexité des connexions entre les neurones est à la base de phénomènes tels que la conscience et la perception. 
	
	\hspace*{1em}Un ensemble de neurones forme un tissu nerveux, tel que celui du cerveau ou de 
	la moelle épinière. Dans certaines parties du corps, les prolongements cellulaires 
	de nombreux neurones sont regroupés avec du tissu conjonctif ; 
	ces prolongements forment un nerf, qui transporte les signaux de nombreux neurones 
	entre le système nerveux et d'autres parties du corps.	
	   
	\hspace*{1em}Les neurones varient en taille et en forme, mais tous partagent des caractéristiques qui permettent la communication entre cellules. Les neurones possèdent des prolongements appelés processus qui assurent les fonctions d'entrée et de sortie des signaux. Ces processus comprennent les dendrites et les axones.

	\hspace*{1em}Le \textbf{corps cellulaire} (ou soma) contient le noyau et les ribosomes, nécessaires à la synthèse des protéines. Les \textbf{dendrites} sont des prolongements ramifiés du corps cellulaire. Dans le système nerveux périphérique (SNP), elles reçoivent des informations sensorielles et les transmettent aux neurones sensoriels. Dans le système nerveux central (SNC), les dendrites et le corps cellulaire reçoivent principalement les signaux des autres neurones. Les dendrites ont une surface étendue grâce à des épines dendritiques, ce qui permet une meilleure réception des signaux et pourrait jouer un rôle dans l'apprentissage et la mémoire.
	
	\hspace*{1em}L'\textbf{axone}, aussi appelé fibre nerveuse, est un long prolongement qui porte les signaux sortants vers les cellules cibles. L'endroit où l'axone commence à partir du corps cellulaire s'appelle le \textbf{segment initial} (ou cône axonique), qui est le point où les signaux électriques sont générés. 
	
	\hspace*{1em}L'axone peut se diviser en branches appelées \textbf{collatérales}. Plus l'axone et ses collatérales sont ramifiés, plus le neurone peut influencer d'autres cellules.
	
\begin{figure}[h]
	\centering
	\includegraphics[width=0.9\linewidth]{pictures/anatomie de neuron 3.png}    
	\caption{Transport axonal le long des microtubules}
\end{figure}    

\begin{figure}[h]
	\centering
	\includegraphics[width=\linewidth]{pictures/anatomie de neuron 1.png}
	\caption{Direction de transmission de l'activité neuronale}
\end{figure}

\begin{figure}[h]
	\centering
	\begin{minipage}{0.7\textwidth}
		\centering
		\includegraphics[width=\linewidth]{pictures/protein 1.png}
	\end{minipage}
	
	\begin{minipage}{0.5\textwidth}
		\centering
		\includegraphics[width=\linewidth]{pictures/protein 2.png}
	\end{minipage}
	\caption{Membrane cytoplasmique du neurone : Protéines membranaires et phospholipides}
\end{figure}

\begin{figure}[h]
	\centering
	\includegraphics[width=\linewidth]{pictures/scale.png}  
	\caption{Échelle de grandeur}
\end{figure}
	  